\section{Referências e rastreabilidade}

\begin{frame}[shrink=4]{Fontes oficiais verificadas nesta revisão}
\small
\begin{itemize}
\item \textbf{ABNT Catálogo}: confirmação da ABNT NBR 5419-1:2026 e consulta das edições indicadas.\\
\href{https://www.abntcatalogo.com.br/}{\texttt{abntcatalogo.com.br}}
\item \textbf{Equatorial Energia Maranhão}: lista vigente das normas de conexão e respectivos anexos, incluindo NT.00001 Rev.09, NT.00002 Rev.10, NT.00004 Rev.08, NT.00020 Rev.06, NT.00030 Rev.03 e NT.00042 Rev.04.\\
\href{https://ma.equatorialenergia.com.br/institucional/leis-e-normas-tecnicas/normas-tecnicas/instalacao-de-padrao/}{\texttt{Equatorial MA — Normas Técnicas de Conexão}}
\item \textbf{Ministério do Trabalho e Emprego}: NR-10 vigente até 31/05/2027 e nova redação com vigência em 01/06/2027.\\
\href{https://www.gov.br/trabalho-e-emprego/pt-br/acesso-a-informacao/participacao-social/conselhos-e-orgaos-colegiados/comissao-tripartite-partitaria-permanente/normas-regulamentadora/normas-regulamentadoras-vigentes/norma-regulamentadora-no-10-nr-10}{\texttt{MTE — página oficial da NR-10}}
\end{itemize}
\begin{caixaAt}
Consulta realizada em 13/08/2026. A lista da distribuidora pode mudar; no início de cada projeto, conferir novamente revisão, anexos e período de transição. Este material resume critérios e não reproduz nem substitui o texto integral das normas ABNT.
\end{caixaAt}
\end{frame}

\begin{frame}{Referências para a loja de shopping}
\small
\begin{itemize}
\item \textbf{ABNT NBR ISO/CIE 8995-1}: iluminação de ambientes de trabalho interiores; confirmar edição vigente no \href{https://www.abntcatalogo.com.br/}{Catálogo ABNT}.
\item \textbf{ABNT NBR 10898}: sistema de iluminação de emergência; usar em conjunto com o projeto de segurança contra incêndio aprovado.
\item \textbf{CBMMA}: \href{https://cbm.ssp.ma.gov.br/unidades-bm/unidades-administrativas/dat/legislacao-tecnica-dat/}{legislação técnica oficial}, incluindo NT 18 para iluminação de emergência e demais normas aplicáveis ao shopping.
\item \textbf{Manual técnico do empreendimento}: ponto de entrega interno, demanda concedida, proteção, medição, materiais, procedimentos de obra e comissionamento.
\item \textbf{Dados fotométricos e fichas dos fabricantes}: arquivos IES/LDT, fluxo, distribuição, potência, IRC, CCT, compatibilidade de controles, torque e manutenção.
\end{itemize}
\begin{caixaAt}
As metas numéricas do estudo da loja são premissas didáticas declaradas. O executivo deve confrontá-las com as edições normativas vigentes, a atividade visual real, a marca e as exigências do shopping.
\end{caixaAt}
\end{frame}
